\documentclass[11pt]{article}
\usepackage[margin=25mm]{geometry}
\usepackage{fontspec,unicode-math}
\setmainfont{Latin Modern Roman}
\setmathfont{Latin Modern Math}
\usepackage{amsmath,amssymb,mathtools}
\usepackage{xurl,hyperref}
\hypersetup{colorlinks=true,urlcolor=blue,linkcolor=blue}
\setcounter{secnumdepth}{0}
\setlength{\parindent}{0pt}
\setlength{\parskip}{5pt}
\setlength{\emergencystretch}{3em}
\newcommand{\sref}[2]{\hyperlink{source:#1:#2}{[S#2]}}
\newcommand{\cataloguescope}[1]{\paragraph{Source question.}#1}
\newcommand{\floor}[1]{\left\lfloor#1\right\rfloor}
\begin{document}
\section{Information–Computation Gaps in Sparse Linear Regression}\label{statement}

\subsection{Definitions and mathematical statement}

\paragraph{Shared notation.}Standard number systems are denoted by $\mathbb N,\mathbb Z,\mathbb Q,\mathbb R,\mathbb C$; index ranges, computational models and specialized notation follow the cited problem source and the definitions below.


Sparse regression models $y=X\beta_*+\varepsilon$ with few nonzero coordinates in $\beta_*$. An information threshold is the least sample scale admitting statistically successful recovery without runtime restrictions; an algorithmic threshold concerns the specified efficient algorithm class.

A learning problem specifies a data distribution, a hypothesis class and a loss. Statistical consistency means that the estimator approaches its population target as the sample size increases; statistical and computational requirements are distinct.

For a discrete random variable $X$, entropy is $H(X)=-\sum_xp(x)\log p(x)$, and mutual information is $I(X;Y)=H(X)+H(Y)-H(X,Y)$. Coding rates, capacities, and exponents depend on the stated channel or source model and error criterion.
\paragraph{Statement recorded in the supplied source.}
In sparse high-dimensional linear regression, when the sample size $n$ lies between the information-theoretic threshold $n_*$ and the algorithmic threshold $n_{\mathrm{alg}}$, is the signal $\beta_*$ information-theoretically recoverable, and is efficient recovery possible? \sref{statement}{2}
\paragraph{Scope qualification.}
The design and noise distributions, sparsity scaling, recovery loss, and threshold definitions require the source.
\paragraph{Source status.}
The supplied catalogue labels this record partially solved. Its source-status metadata was checked against the supplied snapshot on 2026-10-11. This import does not independently establish current open status or mathematical correctness. The catalogue literature assessment is dated 2026-08-21.
\subsection{Short English statement}
Determine whether statistically possible sparse recovery between the two stated thresholds can be performed efficiently.
\subsection{Sources}
\begin{itemize}
\item[\hypertarget{source:statement:1}{S1}] Ulam AI and the UnsolvedMath Contributors, UnsolvedMath: A Curated Collection of Open Mathematics Problems, catalogue entry. Catalogue curation and metadata: CC BY 4.0. The supplied source question is retained; definitions and the English summary are editorial additions.
\url{https://huggingface.co/datasets/ulamai/UnsolvedMath}.
\item[\hypertarget{source:statement:2}{S2}] Network Models: Structure and Function (2018), Oberwolfach Reports, DOI 10.4171/owr/2017/57, pp. 3427--3470
\url{https://doi.org/10.4171/owr/2017/57}.
\item[\hypertarget{source:statement:3}{S3}] High-dimensional regression with binary coefficients
\url{https://arxiv.org/abs/1701.04455}.
\item[\hypertarget{source:statement:4}{S4}] Sparse regression algorithmic barriers
\url{https://arxiv.org/abs/1711.04952}.
\item[\hypertarget{source:statement:5}{S5}] All-or-nothing in sparse linear regression
\url{https://arxiv.org/abs/1903.05046}.
\textit{Underlying Oberwolfach Reports material retains its source-specific licence, including CC BY-SA 4.0 where applicable; see THIRD\_PARTY\_NOTICES.md.}
\end{itemize}
\paragraph{Preparation and review.}
Original question attributed to the cited source and UnsolvedMath contributors. Codex (GPT-6; exact serving revision not exposed) assisted with definitions, English summaries, source inspection and formatting. No new human mathematical authorship or independent proof verification is asserted. Catalogue identity checked against the complete supplied registry and existing repository statements; known flagged repetitions and obvious identical conjectures excluded. Full mathematical-equivalence/literature review has not been performed. Statement preparation and final review followed contribution prompts C and D; automated checks verify structure and identity, not open status.
\label{end:statement}
\end{document}
